\section{Trust and Interaction Experience Metric}
\hsw{We reuse the trust question in xxx.}
\hsw{xxx has been widely used to measure user experience in webpage contexts (cite xxx used it).}
\hsw{However, that scale was designed for webpage contexts. To measure user experience in visualization interaction, we adapted it by changing webpage to interface and aligning wording with visualization tasks.}
\hsw{We then ran fixed-size factor analyses separately for the trust block (4 items) and interaction-experience block (9 items), using 1-factor, 2-factor, and 4-factor solutions in Exp1 (1005 rows; 335 participants), Exp2 (783 rows; 87 participants), and Exp3 (840 rows; 140 participants). We report loadings with $|loading| \ge 0.30$ in \texttt{factor.tex}.}
\hsw{For trust, the 1-factor solution showed consistent high-magnitude loadings across all four items: Exp1 $\{-0.428,\ 0.435,\ -0.918,\ 0.846\}$, Exp2 $\{-0.564,\ 0.471,\ -0.686,\ 0.680\}$, and Exp3 $\{0.494,\ -0.440,\ 0.939,\ -0.726\}$ (same pattern with opposite polarity in Exp3). In 2-factor and 4-factor solutions, extra loadings were sparse and relatively small (mostly $|loading|=0.306$--$0.413$).}
\hsw{For interaction experience, most control/responsiveness/meaningful-interaction items loaded on the first factor across studies (Exp1: $|loading|=0.360$--$0.651$, Exp2: $|loading|=0.301$--$0.440$, Exp3: $|loading|=0.334$--$1.114$ in 1f/2f models). The communication item did not reach the threshold in Exp1 or Exp2, but appeared in Exp3 (2-factor: $0.312$; 4-factor: $-0.341$).}
\hsw{Overall, the data indicate a stable single trust dimension and a mostly single interaction-experience dimension, with a weaker and less stable social-communication component.}
